\documentclass[10pt,a4paper,landscape]{article}
\usepackage[utf8]{inputenc}
\usepackage[T1]{fontenc}
\usepackage{lmodern}
\usepackage[margin=0.7cm]{geometry}
\usepackage{tikz}
\usetikzlibrary{calc,arrows.meta}
\usepackage{parskip}
\usepackage{booktabs}
\pagestyle{empty}
\renewcommand{\familydefault}{\sfdefault}

% ================= MAKRA SYMBOLI (IEC 60617, konwencja EPLAN) =================
\tikzset{every picture/.style={line width=0.55pt}, lbl/.style={font=\tiny, inner sep=0.5pt, fill=white},
tag/.style={font=\scriptsize\bfseries, inner sep=1pt}, zn/.style={font=\tiny, inner sep=0.5pt},
ods/.style={font=\tiny\itshape, inner sep=0.5pt}}

% zestyk zwierny pionowy; dolny zacisk (#1,#2), gorny (#1,#2+1.2)
\newcommand{\ZNO}[2]{\draw (#1,#2) -- (#1,{#2+0.35}) ({#1},{#2+0.35}) -- ({#1-0.3},{#2+1.02}) (#1,{#2+1.02}) -- (#1,{#2+1.2});}
% zestyk rozwierny pionowy
\newcommand{\ZNC}[2]{\draw (#1,#2) -- (#1,{#2+0.35}) ({#1},{#2+0.35}) -- ({#1-0.32},{#2+1.06}) (#1,{#2+0.78}) -- ({#1-0.22},{#2+0.78}) (#1,{#2+0.78}) -- (#1,{#2+1.2});}
% rozlacznik (zwierny + poprzeczka odlacznika na gornym zacisku)
\newcommand{\ROZL}[2]{\ZNO{#1}{#2}\draw ({#1-0.18},{#2+1.02}) -- ({#1+0.18},{#2+1.02});}
% bezpiecznik pionowy
\newcommand{\FUZ}[2]{\draw (#1,#2) -- (#1,{#2+1.2}) ({#1-0.18},{#2+0.22}) rectangle ({#1+0.18},{#2+0.98});}
% bezpiecznik poziomy: lewy zacisk (#1,#2), prawy (#1+1.2,#2)
\newcommand{\FUZH}[2]{\draw (#1,#2) -- ({#1+1.2},#2) ({#1+0.22},{#2-0.18}) rectangle ({#1+0.98},{#2+0.18});}
% zestyk zwierny poziomy: lewy zacisk (#1,#2), prawy (#1+1.2,#2)
\newcommand{\ZNOH}[2]{\draw (#1,#2) -- ({#1+0.18},#2) ({#1+0.18},#2) -- ({#1+0.85},{#2+0.3}) ({#1+0.85},#2) -- ({#1+1.2},#2);}
% cewka pionowa: dolny zacisk (#1,#2), gorny (#1,#2+1.4)
\newcommand{\COILV}[2]{\draw (#1,#2) -- (#1,{#2+0.42}) ({#1-0.42},{#2+0.42}) rectangle ({#1+0.42},{#2+0.98}) (#1,{#2+0.98}) -- (#1,{#2+1.4});}
% przekladnik pradowy: okrag na pionowym przewodzie, centrum (#1,#2)
\newcommand{\CTs}[2]{\draw (#1,#2) circle (0.3); \draw ({#1+0.3},#2) -- ({#1+0.72},#2);}
% zacisk
\newcommand{\TERM}[2]{\draw[fill=white] (#1,#2) circle (0.07);}
% gniazdo (kielich otwarty w gore) + pin z gory
\newcommand{\SOCKup}[2]{\draw ({#1-0.22},{#2+0.1}) arc (180:360:0.22); \draw (#1,{#2+0.55}) -- (#1,{#2+0.06});}
% gniazdo (kielich otwarty w dol) + pin w dol
\newcommand{\SOCKdn}[2]{\draw ({#1-0.22},{#2-0.1}) arc (180:0:0.22); \draw (#1,{#2-0.55}) -- (#1,{#2-0.06});}
% naped E-STOP (grzybek) po lewej zestyku na wys. blade
\newcommand{\ESTOP}[2]{\draw[dashed] ({#1-0.16},{#2+0.7}) -- ({#1-0.6},{#2+0.7}); \draw ({#1-0.6},{#2+0.7}) -- ({#1-0.6},{#2+0.9}) ({#1-0.85},{#2+0.9}) -- ({#1-0.35},{#2+0.9}) ({#1-0.85},{#2+0.9}) arc (180:0:0.25);}
% naped krancowki (rolka)
\newcommand{\ROLL}[2]{\draw[dashed] ({#1-0.16},{#2+0.7}) -- ({#1-0.55},{#2+0.7}); \draw ({#1-0.55},{#2+0.55}) -- ({#1-0.55},{#2+0.85}); \draw ({#1-0.55},{#2+0.55}) circle (0.09);}
% naped termiczny theta
\newcommand{\THETA}[2]{\draw[dashed] ({#1-0.17},{#2+0.7}) -- ({#1-0.5},{#2+0.7}); \node[zn, left] at ({#1-0.5},{#2+0.7}) {$\theta$};}
% przycisk NO (START)
\newcommand{\PUSH}[2]{\draw[dashed] ({#1-0.16},{#2+0.7}) -- ({#1-0.55},{#2+0.7}); \draw ({#1-0.55},{#2+0.55}) -- ({#1-0.55},{#2+0.85}) ({#1-0.8},{#2+0.55}) -- ({#1-0.8},{#2+0.85}) ({#1-0.8},{#2+0.7}) -- ({#1-0.95},{#2+0.7});}
% silnik: kolo M, centrum, podpis
\newcommand{\MOT}[3]{\draw (#1,#2) circle (0.42); \node[font=\scriptsize] at (#1,{#2+0.08}) {M}; \node[font=\tiny] at (#1,{#2-0.16}) {#3};}
% strzalka potencjalu w prawo
\newcommand{\POTR}[2]{\draw (#1,#2) -- ({#1+0.35},{#2+0.12}) -- ({#1+0.35},{#2-0.12}) -- cycle;}

% ================= RAMA ARKUSZA EPLAN =================
% #1 oznaczenie strony, #2 opis strony, #3 nr strony, #4 liczba stron
\newcommand{\ramkaE}[4]{%
\draw[line width=0.8pt] (0,0) rectangle (28.0,18.8);
\draw[line width=0.5pt] (0.5,0.5) rectangle (27.5,18.3);
\foreach \i in {1,...,10}{
  \draw ({0.5+(\i-1)*2.7},18.3) -- ({0.5+(\i-1)*2.7},18.8) ({0.5+(\i-1)*2.7},0.0) -- ({0.5+(\i-1)*2.7},0.5);
  \node[font=\scriptsize] at ({0.5+(\i-1)*2.7+1.35},18.55) {\i};
  \node[font=\scriptsize] at ({0.5+(\i-1)*2.7+1.35},0.25) {\i};}
\foreach \i/\l in {0/A,1/B,2/C,3/D,4/E,5/F}{
  \draw (0.0,{18.3-\i*2.966}) -- (0.5,{18.3-\i*2.966}) (27.5,{18.3-\i*2.966}) -- (28.0,{18.3-\i*2.966});
  \node[font=\scriptsize] at (0.25,{18.3-\i*2.966-1.48}) {\l};
  \node[font=\scriptsize] at (27.75,{18.3-\i*2.966-1.48}) {\l};}
\draw[line width=0.5pt] (14.0,0.5) rectangle (27.5,2.6);
\draw (14.0,1.9) -- (27.5,1.9) (14.0,1.2) -- (23.2,1.2) (18.6,0.5) -- (18.6,2.6) (23.2,0.5) -- (23.2,2.6) (25.6,1.2) -- (25.6,1.9);
\node[font=\tiny, anchor=north west] at (14.05,2.55) {firma};
\node[font=\small\bfseries] at (16.3,2.1) {AMPERE POINT};
\node[font=\tiny, anchor=north west] at (14.05,1.85) {projekt};
\node[font=\scriptsize] at (16.3,1.45) {AMPERE\_POINT\_custom\_device};
\node[font=\tiny, anchor=north west] at (14.05,1.15) {nr projektu};
\node[font=\scriptsize] at (16.3,0.75) {AP-CD-001};
\node[font=\tiny, anchor=north west] at (18.65,2.55) {opis strony};
\node[font=\scriptsize, align=left, anchor=west] at (18.75,2.1) {#2};
\node[font=\tiny, anchor=north west] at (18.65,1.85) {data / edytor};
\node[font=\scriptsize, anchor=west] at (18.75,1.45) {09.07.2026 \; / \; D.~Cekała};
\node[font=\tiny, anchor=north west] at (18.65,1.15) {stan: wersja 5 (wariant budżetowy v6)};
\node[font=\tiny, anchor=north west] at (23.25,2.55) {oznaczenie strony};
\node[font=\small\bfseries] at (24.4,2.1) {#1};
\node[font=\tiny, anchor=north west] at (23.25,1.85) {norma};
\node[font=\scriptsize] at (24.4,1.45) {IEC 60617};
\node[font=\tiny, anchor=north west] at (25.65,1.85) {strona};
\node[font=\scriptsize] at (26.55,1.45) {#3 / #4};
\node[font=\scriptsize\bfseries, anchor=west] at (23.3,0.85) {custom device --- stacja testowa EVSE};
}

\begin{document}

% ============ STRONA TYTULOWA ============
\begin{center}
\vspace*{1.4cm}
{\LARGE\bfseries AMPERE POINT --- custom device}\\[6pt]
{\Large Schemat elektryczny --- wersja 5 (styl EPLAN, wariant budżetowy)}\\[10pt]
{\normalsize 9 lipca 2026}\\[12pt]
\begin{minipage}{0.88\textwidth}\small
Wersja 5 $=$ wersja 4 dostosowana do koncepcji v6 (pełny plan realizacji, wariant budżetowy
Z16--Z19): prostowniki \textbf{Eltek Flatpack2 48/2000~HE} (używane, CAN), \textbf{pakiet bateryjny
16S LFP 105~Ah składany samodzielnie} (rozbudowa 2P), falownik \textbf{W2: używany UPS on-line
48~V} zasilający wydzielone obwody wyspy, \textbf{pole przełączeń -X6} (ręczne mostki L2$\to$L1 do
testu 1F 32~A) zamiast stycznika -K14, \textbf{bez zrzutu cieplnego} (bramka SOC w firmware).
Arkusze E1/E2 --- bez zmian topologii. Objaśnienie wszystkich aparatów i przewodnik uruchomienia:
\texttt{...schemat\_elektryczny\_v5\_objasnienie\_i\_przewodnik.pdf}.\\[6pt]
\textbf{E1} --- tor mocy 3F: X1 $\to$ Q0 $\to$ F1--F3 $\to$ B1 $\to$ T1--T3 $\to$ K1 $\to$ X2; odczep X3 (7 żył) z -K4 w linii CP.\\
\textbf{E2} --- łańcuch bezpieczeństwa 24 V (termiki -F11.1..8 na radiatorach przetwornic).\\
\textbf{E3} --- szafa magazynu: -K11..-K13 (+pole -X6) $\to$ -PR1..-PR6 (Flatpack2, CAN) $\to$
DC 48 V (-F13, -KB, -Q1) $\to$ -GB1 (16S LFP) $\to$ -INV (UPS) $\to$ obwody wyspy.
\end{minipage}\\[10pt]
{\small\textbf{Zmiany w wersji 5 względem wersji 4}}\\[3pt]
{\scriptsize
\begin{tabular}{clp{11.6cm}}
\toprule
\# & Arkusz & Zmiana\\
\midrule
1 & E3 & typy prostowników: -PR1..-PR6 $=$ 2$\times$ Flatpack2 48/2000 HE na fazę (używane, CAN 125 kb/s)\\
2 & E3 & stycznik -K14 usunięty $\to$ \textbf{-X6 pole przełączeń} (ręczne mostki; konfiguracja weryfikowana pomiarem -U2 przy 25\% mocy)\\
3 & E3 & usunięty zrzut CWU (-K5, -E13) $\to$ bramka SOC w firmware ($\le$40\% przed soakiem 3F / $\le$60\% 1F)\\
4 & E3 & -GB1: pakiet 16S LFP 105 Ah (5,4 kWh; 2P $=$ 10,8 kWh), BMS $\to$ kontaktor -KB i telemetria do -A1\\
5 & E3 & -INV $=$ falownik W2: używany UPS on-line 48 V; wyjście na wydzielone obwody wyspy (przełącznik sieć--0--wyspa)\\
6 & E3 & rząd cewek zredukowany: -K11..-K13, -K30..-K33, -K3 (bez -K14/-K5)\\
\bottomrule
\end{tabular}}
\end{center}
\newpage
% ============ ARKUSZ E1 — TOR MOCY ============
\begin{center}
\begin{tikzpicture}[x=1cm,y=1cm]
\ramkaE{=CA1+EAA/2}{E1 --- tor mocy 3F}{2}{4}
% ---- X1 gniazdo pojazdowe ----
\foreach \x/\n/\z in {5.0/L1/1, 7.2/L2/2, 9.4/L3/3, 11.6/N/4, 13.8/PE/5, 15.4/CP/6, 16.2/PP/7}{
  \draw (\x,17.65) -- (\x,17.25);
  \SOCKup{\x}{16.7}
  \node[zn] at (\x,17.85) {\n};
  \node[zn, right] at ({\x+0.1},16.55) {\z};
  \draw (\x,16.6) -- (\x,16.3);}
\draw[dashed] (4.3,16.25) rectangle (16.9,17.45);
\node[tag, left] at (4.25,17.3) {-X1};
\node[lbl, align=left] at (2.3,16.6) {gniazdo pojazdowe\\ Type 2 32 A 3F\\ (wtyk kabla DUT)};
\node[lbl, align=left] at (2.3,15.5) {blokada wtyku\\ -Y1 $\to$ /E3.3};
% ---- piony ----
\draw (5.0,16.3) -- (5.0,14.8) (7.2,16.3) -- (7.2,14.8) (9.4,16.3) -- (9.4,14.8) (11.6,16.3) -- (11.6,14.8);
\draw (13.8,16.3) -- (13.8,6.0); % PE
\node[lbl] at (14.15,9.0) {PE};
% ---- odczep X3 (7 zyl) ----
\draw[fill] (5.0,16.0) circle (0.06); \draw (5.0,16.0) -- (17.0,16.0);
\draw[fill] (7.2,15.7) circle (0.06); \draw (7.2,15.7) -- (17.0,15.7);
\draw[fill] (9.4,15.4) circle (0.06); \draw (9.4,15.4) -- (17.0,15.4);
\draw[fill] (11.6,15.1) circle (0.06); \draw (11.6,15.1) -- (17.0,15.1);
\draw[fill] (13.8,14.8) circle (0.06); \draw (13.8,14.8) -- (17.0,14.8);
% CP przez zestyk -K4
\draw (15.4,16.3) -- (15.4,15.75);
\ZNO{15.4}{14.55}
\node[tag, left] at (15.05,15.3) {-K4};
\node[ods, left] at (15.05,14.95) {/E2.6};
\draw (15.4,14.55) -- (15.4,14.5) -- (17.0,14.5);
% PP
\draw (16.2,16.3) -- (16.2,14.2) -- (17.0,14.2);
\foreach \y/\z in {16.0/1, 15.7/2, 15.4/3, 15.1/4, 14.8/5, 14.5/6, 14.2/7}{\TERM{17.0}{\y}\node[zn, right] at (17.1,\y) {\z};}
\draw[dashed] (16.75,13.95) rectangle (17.75,16.25);
\node[tag] at (17.25,16.5) {-X3};
\node[lbl, align=left, anchor=west] at (18.0,15.75) {odczep Type 2 $\to$ PM701E (7 żył, $\le$1 m)\\ 1=L1, 2=L2, 3=L3, 4=N, 5=PE, 6=CP, 7=PP};
\node[ods, anchor=west] at (18.0,15.05) {SIGNAL OUTPUT (CP) $\to$ /E3.3};
\node[ods, anchor=west] at (18.0,14.65) {gniazda bananowe PM701E = przyłącze UT595};
\node[ods, anchor=west] at (18.0,14.25) {mostek serwisowy CP (obejście -K4) na listwie -X3};
% ---- Q0 rozlacznik 4P ----
\ROZL{5.0}{13.6}\ROZL{7.2}{13.6}\ROZL{9.4}{13.6}\ROZL{11.6}{13.6}
\draw[dashed] (4.55,14.25) -- (11.45,14.25);
\node[tag, left] at (4.2,14.2) {-Q0};
\node[lbl, left] at (4.2,13.8) {rozłącznik 40 A};
\foreach \x/\a/\b in {5.0/1/2, 7.2/3/4, 9.4/5/6, 11.6/7/8}{
  \node[zn, right] at ({\x+0.08},14.72) {\a}; \node[zn, right] at ({\x+0.08},13.7) {\b};}
% ---- F1-F3 ----
\draw (5.0,13.6) -- (5.0,13.1) (7.2,13.6) -- (7.2,13.1) (9.4,13.6) -- (9.4,13.1);
\draw (11.6,13.6) -- (11.6,11.4);
\FUZ{5.0}{11.9}\FUZ{7.2}{11.9}\FUZ{9.4}{11.9}
\node[tag, left] at (4.55,12.5) {-F1};\node[tag, left] at (6.75,12.5) {-F2};\node[tag, left] at (8.95,12.5) {-F3};
\node[lbl, left] at (4.2,12.05) {gG 32 A};
% ---- B1 okno CDSR ----
\draw (5.0,11.9) -- (5.0,9.9) (7.2,11.9) -- (7.2,9.9) (9.4,11.9) -- (9.4,9.9);
\draw (11.6,11.4) -- (11.6,9.9);
\draw[line width=1.0pt] (4.3,10.6) rectangle (12.3,11.35);
\node[tag, left] at (4.2,11.0) {-B1};
\node[zn] at (8.3,10.97) {CDSR 0.07-TP};
\node[zn, align=left, anchor=south west] at (13.35,11.2) {czujnik różnicowy AC/DC (RCM), okno: L1+L2+L3+N};
\draw[-{Stealth}] (12.3,11.0) -- (13.1,11.0); \node[ods, right] at (13.1,11.0) {SPI /E3.3};
% ---- odczep L1 do generatora A3 (za oknem B1) ----
\draw[fill] (5.0,10.05) circle (0.06);
\draw[-{Stealth}] (5.0,10.05) -- (3.6,10.05);
\node[ods, align=right, anchor=east] at (3.55,10.05) {$\to$ -A3 /E3.2\\ (przez -K30)};
% ---- T1-T3 ----
\draw (5.0,9.9) -- (5.0,8.8) (7.2,9.9) -- (7.2,8.8) (9.4,9.9) -- (9.4,8.8) (11.6,9.9) -- (11.6,8.8);
\CTs{5.0}{9.35}\CTs{7.2}{9.35}\CTs{9.4}{9.35}
\node[tag, left] at (4.55,9.35) {-T1};\node[tag, left] at (6.75,9.35) {-T2};\node[tag, left] at (8.95,9.35) {-T3};
\node[ods] at (6.1,8.85) {/E3.3};\node[ods] at (8.3,8.85) {/E3.3};\node[ods] at (10.5,8.85) {/E3.3};
\node[lbl, left] at (4.2,9.0) {40 A / 1 V};
% ---- K1 stycznik glowny ----
\ZNO{5.0}{7.6}\ZNO{7.2}{7.6}\ZNO{9.4}{7.6}\ZNO{11.6}{7.6}
\draw[dashed] (4.55,8.25) -- (11.45,8.25);
\node[tag, left] at (4.2,8.35) {-K1};
\node[ods, left] at (4.2,7.95) {/E2.5};
\foreach \x/\a/\b in {5.0/1/2, 7.2/3/4, 9.4/5/6, 11.6/7/8}{
  \node[zn, right] at ({\x+0.08},8.72) {\a}; \node[zn, right] at ({\x+0.08},7.7) {\b};}
% ---- zejscie do X2 ----
\draw (5.0,7.6) -- (5.0,6.0) (7.2,7.6) -- (7.2,6.0) (9.4,7.6) -- (9.4,6.0) (11.6,7.6) -- (11.6,6.0);
\foreach \x/\n/\z in {5.0/L1/1, 7.2/L2/2, 9.4/L3/3, 11.6/N/4, 13.8/PE/5}{
  \SOCKdn{\x}{5.9}
  \draw (\x,5.35) -- (\x,5.0);
  \node[zn, right] at ({\x+0.1},6.05) {\z};
  \node[zn] at (\x,4.82) {\n};}
\draw[dashed] (4.3,5.15) rectangle (14.6,6.3);
\node[tag, left] at (4.25,6.15) {-X2};
\node[lbl, align=left] at (2.4,5.7) {gniazdo CEE 32 A 5p\\ $\to$ moduł B /E3.3};
\node[lbl] at (8.3,4.4) {przewód 5$\times$6 mm$^2$ H07RN-F};
\node[lbl, left] at (4.6,7.0) {6 mm$^2$};
% ---- uwagi ----
\node[lbl, align=left, anchor=north west] at (18.6,13.6) {\textbf{Uwagi:}\\ N bez zabezpieczenia --- rozłączany\\ czterobiegunowo (Q0, K1)\\ PE nieprzerwane od X1 do X2\\ B1: pomiar RCM, próg awaryjny $\to$ /E3.3\\ T1--T3: pomiar prądów $\to$ U2 /E3.3\\ K1: cewka i lustro styków $\to$ /E2.5};
\node[lbl, align=left, anchor=north west] at (18.6,10.7) {\textbf{Ostrzeżenie (v3):} -Q0 nie odłącza odczepu\\ -X3/PM701E --- gniazda bananowe mogą być\\ pod napięciem w stanie C niezależnie od\\ -Q0 i -K1 $\to$ tabliczka ostrzegawcza na panelu.};
\node[lbl, align=left, anchor=north west] at (18.6,8.9) {\textbf{-K4 (v3):} zestyk w linii CP odczepu;\\ E-STOP / zanik 24 V $\to$ przerwanie CP $\to$\\ DUT widzi stan A i otwiera swój stycznik\\ (magistrala beznapięciowa). Cewka: /E2.6.};
\end{tikzpicture}
\end{center}
\newpage
% ============ ARKUSZ E2 — LANCUCH BEZPIECZENSTWA ============
\begin{center}
\begin{tikzpicture}[x=1cm,y=1cm]
\ramkaE{=CA1+EAA/3}{E2 --- łańcuch bezp. 24 V}{3}{4}
% ---- szyny ----
\draw (1.4,16.9) -- (16.0,16.9); \node[tag, left] at (1.35,16.9) {+24 V};
\POTR{16.0}{16.9} \node[ods, right] at (16.4,16.9) {z G1 /E3.5};
\draw (1.4,3.6) -- (16.0,3.6); \node[tag, left] at (1.35,3.6) {0 V};
\POTR{16.0}{3.6} \node[ods, right] at (16.4,3.6) {G1 /E3.5};
% ================= sciezka 1 (kol.1): S0 -> termiki -> S1 =================
\draw[fill] (3.0,16.9) circle (0.06);
\draw (3.0,16.9) -- (3.0,16.5);
\ZNC{3.0}{15.3}\ESTOP{3.0}{15.3}
\node[tag, right] at (3.2,15.9) {-S0};
\node[zn, right] at (3.08,16.35) {11}; \node[zn, right] at (3.08,15.4) {12};
\node[lbl, right] at (3.2,15.55) {E-STOP (panel A)};
\draw (3.0,15.3) -- (3.0,14.8);
\ZNC{3.0}{13.6}\THETA{3.0}{13.6}
\node[tag, right] at (3.2,14.2) {-F11.1};
\draw (3.0,13.6) -- (3.0,13.35);
\node[font=\scriptsize] at (3.0,13.05) {$\vdots$};
\draw (3.0,12.75) -- (3.0,12.7);
\ZNC{3.0}{11.5}\THETA{3.0}{11.5}
\node[tag, right] at (3.2,12.1) {-F11.12};
\draw[dashed] (1.95,11.25) rectangle (4.75,14.95);
\node[lbl, align=left, anchor=west] at (4.95,12.35) {termiki KSD301 radiatorów\\ przetwornic (NC szeregowo);\\ moduł B, złącze -X5 /E3.6};
\draw (3.0,11.5) -- (3.0,10.9);
\ZNC{3.0}{9.7}\ROLL{3.0}{9.7}
\node[tag, right] at (3.2,10.3) {-S1};
\node[lbl, right] at (3.2,9.95) {krańcówka pokrywy};
\node[zn, right] at (3.08,10.75) {11}; \node[zn, right] at (3.08,9.8) {12};
\draw (3.0,9.7) -- (3.0,8.9);
\draw[fill] (3.0,8.9) -- (2.85,8.55) -- (3.15,8.55) -- cycle;
\node[ods] at (3.0,8.25) {/E2.3};
% ================= sciezka 3 (kol.3): K2 -> K10 -> A2 -> [S2||RB] -> cewka RB =================
\draw[fill] (8.4,16.4) -- (8.25,16.05) -- (8.55,16.05) -- cycle;
\node[ods] at (8.4,16.6) {/E2.1};
\draw (8.4,16.05) -- (8.4,15.6);
\ZNO{8.4}{14.4}
\node[tag, left] at (8.1,15.0) {-K2};
\node[ods, left] at (8.1,14.6) {/E3.8};
\node[lbl, right] at (8.6,15.0) {potwierdzenie wentylatorów};
\node[zn, right] at (8.48,15.5) {13}; \node[zn, right] at (8.48,14.5) {14};
\draw (8.4,14.4) -- (8.4,13.9);
\ZNO{8.4}{12.7}
\node[tag, left] at (8.1,13.3) {-K10};
\node[ods, left] at (8.1,12.9) {/E3.5};
\node[lbl, right] at (8.6,13.3) {zezwolenie MCU};
\node[zn, right] at (8.48,13.8) {13}; \node[zn, right] at (8.48,12.8) {14};
\draw (8.4,12.7) -- (8.4,12.2);
\ZNO{8.4}{11.0}
\node[tag, left] at (8.1,11.6) {-A2};
\node[ods, left] at (8.1,11.2) {/E3.5};
\node[lbl, right] at (8.6,11.6) {watchdog impulsowy};
\draw (8.4,11.0) -- (8.4,10.5);
% galaz rownolegla
\draw[fill] (8.4,10.5) circle (0.06);
\draw (8.4,10.5) -- (10.6,10.5);
\PUSH{8.4}{8.9}\ZNO{8.4}{8.9}
\node[tag, anchor=east] at (7.3,10.15) {-S2};
\node[lbl, anchor=east] at (7.3,9.8) {START};
\node[zn, right] at (8.48,10.0) {13}; \node[zn, right] at (8.48,9.0) {14};
\draw (10.6,10.5) -- (10.6,10.1);
\ZNO{10.6}{8.9}
\node[tag, right] at (10.8,9.5) {-RB};
\node[ods, right] at (10.8,9.15) {13-14 \; /E2.4};
\draw (10.6,8.9) -- (10.6,8.5) -- (8.4,8.5);
\draw[fill] (8.4,8.5) circle (0.06);
\node[lbl, right] at (10.8,8.7) {samopodtrzymanie};
\draw (8.4,8.9) -- (8.4,8.0);
\COILV{8.4}{6.6}
\node[tag, left] at (7.9,7.3) {-RB};
\node[zn, right] at (8.48,7.85) {A1}; \node[zn, right] at (8.48,6.75) {A2};
\node[lbl, left] at (7.9,6.9) {przekaźnik bezpieczeństwa};
\draw (8.4,6.6) -- (8.4,3.6);
\draw[fill] (8.4,3.6) circle (0.06);
% lustro stykow RB
\node[ods, align=left, anchor=north west] at (8.7,6.1) {\textbf{-RB}\\ 13-14 \; /E2.4\\ 23-24 \; /E2.5};
% ================= sciezka 5 (kol.5): RB:23-24 -> cewka K1 =================
\draw[fill] (13.0,16.9) circle (0.06);
\draw (13.0,16.9) -- (13.0,15.8);
\ZNO{13.0}{14.6}
\node[tag, left] at (12.7,15.2) {-RB};
\node[ods, left] at (12.7,14.8) {/E2.3};
\node[zn, right] at (13.08,15.7) {23}; \node[zn, right] at (13.08,14.7) {24};
\node[lbl, right] at (13.2,15.15) {styk wymuszony};
\draw (13.0,14.6) -- (13.0,13.2);
\COILV{13.0}{11.8}
\node[tag, left] at (12.5,12.5) {-K1};
\node[zn, right] at (13.08,13.05) {A1}; \node[zn, right] at (13.08,11.95) {A2};
\node[lbl, right] at (13.5,12.5) {stycznik główny /E1};
\draw (13.0,11.8) -- (13.0,3.6);
\draw[fill] (13.0,3.6) circle (0.06);
% lustro stykow K1
\node[ods, align=left, anchor=north west] at (13.3,10.9) {\textbf{-K1}\\ 1-2 \; /E1.2\\ 3-4 \; /E1.3\\ 5-6 \; /E1.4\\ 7-8 \; /E1.5};
% ================= sciezka 6 (kol.6): S0:21-22 -> cewka K4 (v3) =================
\draw[fill] (15.6,16.9) circle (0.06);
\draw (15.6,16.9) -- (15.6,15.8);
\ZNC{15.6}{14.6}
\node[tag, left] at (15.3,15.2) {-S0};
\node[ods, left] at (15.3,14.85) {/E2.1};
\node[zn, right] at (15.68,15.7) {21}; \node[zn, right] at (15.68,14.7) {22};
\node[lbl, right] at (15.85,15.45) {E-STOP, tor 2};
\draw (15.6,14.6) -- (15.6,12.4);
\COILV{15.6}{11.0}
\node[tag, left] at (15.1,11.7) {-K4};
\node[zn, right] at (15.68,12.25) {A1}; \node[zn, right] at (15.68,11.15) {A2};
\node[lbl, right] at (16.1,11.7) {przerwanie CP (v3)};
\draw (15.6,11.0) -- (15.6,3.6);
\draw[fill] (15.6,3.6) circle (0.06);
% lustro stykow K4
\node[ods, align=left, anchor=north west] at (15.9,10.6) {\textbf{-K4}\\ styk CP \; /E1.6};
% ---- adnotacje ----
\node[lbl, align=left, anchor=north west] at (17.6,14.0) {\textbf{Zasada:} tor 24 V cewki -K1 przechodzi wyłącznie\\ przez styki sprzętowe. Awaria członu, zanik 24 V,\\ zawieszenie firmware (zanik impulsów -A2)\\ lub E-STOP $\Rightarrow$ -RB odpada $\Rightarrow$ -K1 otwiera\\ tor mocy (/E1). Ponowne załączenie tylko -S2.};
\node[lbl, align=left, anchor=north west] at (17.6,11.1) {\textbf{-K4 (v3):} cewka wprost z 24 V przez -S0:21-22\\ (poza łańcuchem RB) --- podtrzymuje ciągłość CP;\\ E-STOP lub zanik 24 V $\Rightarrow$ przerwanie CP /E1.6\\ $\Rightarrow$ DUT sam zdejmuje napięcie z magistrali.};
\node[lbl, align=left, anchor=north west] at (17.6,8.6) {\textbf{Odczyty diagnostyczne (optoizolacja):}\\ stan -S0, -F11.x, -S1, -K2 $\to$ A1 /E3.4\\ (nie ingerują w tor łańcucha)};
\node[lbl, align=left, anchor=north west] at (17.6,7.0) {\textbf{Poza łańcuchem:} B1 (RCM) $\to$ zdjęcie -K10 /E3.5;\\ bezpieczniki gG /E1.2; rozłącznik -Q0 /E1.2;\\ blokada -Y1 wtyku Type 2 w stanie C /E3.3};
\end{tikzpicture}
\end{center}
\newpage
% ============ ARKUSZ E3 — SZAFA REKUPERACJI + STEROWANIE ============
\begin{center}
\begin{tikzpicture}[x=1cm,y=1cm]
\ramkaE{=CA1+EAA/4}{E3 --- magazyn energii i sterowanie}{4}{4}
\draw[dashed] (7.5,8.1) rectangle (27.0,17.7);
\node[tag, anchor=south west] at (7.6,17.75) {MODUŁ B --- szafa magazynu energii, wariant budżetowy (zasilanie z -X2 /E1.2--5, sterowanie -X5, CAN)};
\foreach \x/\f/\k/\p in {9.2/{L1 /E1.2}/-K11/{-PR1/-PR2}, 12.6/{L2 /E1.3}/-K12/{-PR3/-PR4}, 16.0/{L3 /E1.4}/-K13/{-PR5/-PR6}}{
  \node[ods] at (\x,17.35) {\f};
  \draw (\x,17.2) -- (\x,16.6);
  \ZNO{\x}{15.4}
  \node[tag, left] at ({\x-0.35},16.05) {\k};
  \node[ods, left] at ({\x-0.35},15.65) {/E3.6};
  \draw (\x,15.4) -- (\x,15.0);
  \draw[fill=white] ({\x-0.16},15.06) rectangle ({\x+0.16},15.34);
  \draw ({\x-1.5},13.3) rectangle ({\x+1.5},15.0);
  \node[zn] at (\x,14.65) {\p};
  \node[zn] at (\x,14.3) {2$\times$ Flatpack2 48/2000 HE};
  \node[zn] at (\x,13.95) {używane, PFC, CAN 125 kb/s};
  \draw ({\x-0.9},13.3) -- ({\x-0.9},12.85);
  \draw[fill] ({\x-0.9},12.85) circle (0.05);
  \draw ({\x+0.9},13.3) -- ({\x+0.9},11.9);
  \draw[fill] ({\x+0.9},11.9) circle (0.05);
  \draw ({\x+0.3},13.3) -- ({\x+0.3},11.45);
  \draw[fill] ({\x+0.3},11.45) circle (0.05);}
\node[tag, left] at (8.85,15.2) {-X6};
\node[ods, align=left, anchor=west] at (17.75,15.35) {-X6: pole przełączeń --- mostki L2$\to$L1\\ dla 1F 32 A (kontrola: pomiar -U2)};
\POTR{7.7}{12.85}\draw (8.05,12.85) -- (17.6,12.85); \node[ods, left] at (7.65,12.85) {N /E1.5};
\draw (8.2,11.9) -- (22.8,11.9); \node[zn, left] at (8.15,11.9) {+48 V};
\draw (8.2,11.45) -- (22.6,11.45); \node[zn, left] at (8.15,11.45) {0 V};
\draw[fill] (18.0,11.9) circle (0.05);
\draw (18.0,11.9) -- (18.0,11.0);
\FUZH{18.0}{11.0}
\node[tag] at (18.55,11.5) {-F13.1};
\ZNOH{19.3}{11.0}
\node[tag] at (19.9,11.35) {-KB};
\node[ods] at (19.9,10.7) {kontaktor BMS};
\draw (20.5,11.0) -- (21.2,11.0);
\ROZL{21.2}{9.8}
\node[tag, right] at (21.4,10.5) {-Q1};
\node[ods, right] at (21.4,10.15) {rozłącznik DC};
\draw (19.8,8.4) rectangle (24.2,9.8);
\node[zn] at (22.0,9.45) {\textbf{-GB1}: pakiet 16S LFP 105 Ah (DIY)};
\node[zn] at (22.0,9.1) {5,4 kWh (2P: 10,8); BMS $\to$ -A1 i -KB};
\node[zn] at (22.0,8.72) {temperatury: -B5.x (1-Wire)};
\draw[fill] (19.0,11.45) circle (0.05);
\draw (19.0,11.45) -- (19.0,10.8);
\FUZ{19.0}{9.6}
\node[tag, left] at (18.8,10.2) {-F13.2};
\draw (19.0,9.6) -- (19.0,8.15) -- (20.9,8.15) -- (20.9,8.4);
\draw[fill] (22.8,11.9) circle (0.05);
\draw (22.8,11.9) -- (22.8,13.4) -- (23.2,13.4);
\draw[fill] (22.6,11.45) circle (0.05);
\draw (22.6,11.45) -- (22.6,12.9) -- (23.2,12.9);
\draw (23.2,12.6) rectangle (26.8,15.8);
\node[zn] at (25.0,15.45) {\textbf{-INV} --- falownik W2: UPS on-line 48 V};
\node[zn] at (25.0,15.1) {używany, 5--6 kVA; DC wprost z -GB1};
\node[zn] at (25.0,14.75) {wyjście: wydzielone obwody wyspy};
\draw[-{Stealth}] (26.2,15.8) -- (26.2,16.7);
\node[ods, right] at (25.0,16.95) {$\to$ obwody wyspy (sieć--0--wyspa)};
\node[ods, align=left, anchor=north west] at (8.1,11.15) {bez zrzutu cieplnego: bramka SOC w firmware\\ ($\le$40\% przed soakiem 3F; 1F: $\le$60\%)};
\MOT{18.2}{16.9}{1$\sim$}
\node[ods, anchor=west] at (18.75,17.1) {wentylatory strat -M3/-M4 ($\sim$0,5--1,2 kW ciepła):};
\node[ods, anchor=west] at (18.75,16.7) {X4:L $\to$ -F5 $\to$ styk -K3 /E3.8 $\to$ czujnik -K2 (I$>$) $\to$ /E2.3};
\node[ods, anchor=west] at (18.75,16.3) {termiki $\theta$ -F11.1..8 na radiatorach -PR1..6 i -INV $\to$ /E2.1 (-X5)};
\node[ods, left] at (1.3,10.0) {X4:1 (L)};
\node[ods, left] at (1.3,9.6) {X4:2 (N)};
\node[ods, left] at (1.3,9.2) {X4:3 (PE)};
\foreach \y in {10.0,9.6,9.2}{\TERM{1.45}{\y}}
\draw (1.52,10.0) -- (1.6,10.0);
\FUZH{1.6}{10.0}
\node[tag] at (2.2,10.35) {-F4};
\node[zn] at (2.2,9.6) {T2A};
\draw (1.52,9.6) -- (2.9,9.6);
\draw (1.52,9.2) -- (2.9,9.2);
\draw (2.8,10.0) -- (2.9,10.0);
\draw (2.9,8.9) rectangle (5.5,10.3);
\node[font=\scriptsize\bfseries, align=center] at (4.2,9.85) {-G1};
\node[zn, align=center] at (4.2,9.35) {zasilacz\\ 24 / 5 / 3,3 V};
\draw[-{Stealth}] (5.5,9.95) -- (6.1,9.95); \node[ods, right, align=left] at (6.1,10.1) {24 V $\to$ /E2.1,\\ cewki, -Y1, -K4};
\draw[-{Stealth}] (5.5,9.25) -- (6.1,9.25); \node[ods, right] at (6.1,9.25) {5 / 3,3 V $\to$ -A1};
\node[lbl, align=left, anchor=north west] at (2.9,8.6) {zasilanie logiki z przewodu Schuko\\ (nie z badanej ładowarki)};
\draw (1.0,11.4) rectangle (6.4,13.1);
\node[font=\scriptsize\bfseries] at (3.7,12.75) {-A3 --- generator upływu};
\node[zn, align=left, anchor=north west] at (1.15,12.5) {DC 0--20 mA / AC (rampa z DAC);\\ pomiar: bocznik + AMC1301 $\to$ -A1;\\ przyłącza: L1 za -B1 (/E1.2, przez -K30), PE};
\node[ods, anchor=north west] at (1.15,11.35) {test RDC-DD 6 mA DC i RCD ładowarki};
\draw (1.5,3.4) rectangle (7.0,7.6);
\node[font=\scriptsize\bfseries] at (4.25,7.3) {-A1 \; ESP32-S3 --- płyta logiki};
\foreach \y/\p/\t in {6.85/SPI/{-B1 /E1.3; -U2 ATM90E36A $\leftarrow$ -T1\ldots-T3 /E1.3},
                      6.35/CAN/{TWAI izol.: -PR1..-PR6 (Flatpack2) + BMS /E3.4},
                      5.85/I2C/{-B2 MLX90640; -B3, -B4 MLX90614; -U3 MCP23017},
                      5.35/UART/{-P1 HMI DWIN 7'' (START, LED, buzzer)},
                      4.85/{1-W}/{-B5.1\ldots8 DS18B20 (szyny, -K1, bateria, otoczenie)},
                      4.35/ADC/{CP, PP z -X3 /E1.6 (dzielnik + ograniczniki)},
                      3.85/DI/{odczyty łańcucha (opto) /E2.1: -S0, -F11.x, -S1, -K2}}{
  \TERM{7.0}{\y}
  \node[zn, left] at (6.9,\y) {\p};
  \draw (7.07,\y) -- (7.6,\y);
  \node[ods, right, align=left] at (7.6,\y) {\t};}
\draw[densely dashed, -{Stealth}] (9.2,13.3) -- (9.2,12.55) node[ods, below, align=center] {CAN 125 kb/s\\ $\to$ -A1 /E3.3};
\node[zn] at (3.2,3.55) {DO1};
\node[zn] at (4.0,3.55) {DO2};
\node[zn] at (4.8,3.55) {DO3};
\TERM{3.2}{3.4}\TERM{4.0}{3.4}\TERM{4.8}{3.4}
\draw (3.2,3.33) -- (3.2,2.9) -- (12.4,2.9) -- (12.4,3.3);
\draw (4.0,3.33) -- (4.0,2.55) -- (13.65,2.55) -- (13.65,3.3);
\draw (4.8,3.33) -- (4.8,2.2) -- (8.2,2.2);
\COILV{12.4}{3.3}
\node[tag, right] at (12.9,4.35) {-K10};
\node[ods, right] at (12.9,4.0) {styk /E2.3};
\draw (12.4,4.7) -- (12.4,4.95); \node[ods, above] at (12.4,4.95) {24 V};
\draw (13.25,3.3) rectangle (14.05,4.1); \node[zn, align=center] at (13.65,3.7) {-A2\\ WDT};
\node[ods, anchor=east] at (13.2,2.95) {styk /E2.3};
\draw (13.65,4.1) -- (13.65,4.35); \node[ods, anchor=east] at (13.55,4.55) {24 V};
\draw[fill] (8.2,2.2) circle (0.05);
\draw (8.2,2.2) -- (8.2,2.01) ({8.2-0.42},1.45) rectangle ({8.2+0.42},2.01) (8.2,1.45) -- (8.2,1.2) node[zn, below]{24 V};
\node[tag, right] at (8.75,1.85) {-Y1};
\node[ods, right] at (8.75,1.5) {blokada wtyku -X1 /E1.6 --- 24 V przez DO3};
\node[lbl, align=left, anchor=west] at (1.6,1.55) {DO1--DO3 (opto, low-side):\\ zezwolenie / watchdog / blokada};
\draw (14.4,6.4) -- (21.2,6.4); \node[ods, left] at (14.35,6.4) {24 V};
\foreach \i/\x in {11/14.9, 12/15.6, 13/16.3, 30/17.5, 31/18.2, 32/18.9, 33/19.6, 3/20.5}{
  \draw[fill] (\x,6.4) circle (0.05);
  \draw (\x,6.4) -- (\x,6.15);
  \draw ({\x-0.28},5.55) rectangle ({\x+0.28},6.15);
  \draw (\x,5.55) -- (\x,5.3);
  \node[zn, rotate=90, anchor=east] at ({\x-0.16},5.25) {-K\i};}
\draw (14.4,4.35) -- (21.2,4.35); \node[ods, left] at (14.35,4.35) {0 V};
\foreach \x in {14.9,15.6,16.3,17.5,18.2,18.9,19.6,20.5}{
  \draw (\x,5.3) -- (\x,4.35); \draw[fill] (\x,4.35) circle (0.05);}
\draw (13.0,6.9) rectangle (17.4,7.7); \node[zn, align=center] at (15.2,7.3) {-U5 ULN2803 $\leftarrow$ -U3 (I2C -A1)};
\draw[-{Stealth}] (15.2,6.9) -- (15.2,6.55);
\node[lbl, align=left, anchor=north west] at (14.4,3.95) {styki: -K11..-K13 /E3.4--6 (sekcje faz); -X6 pole przełączeń 1F 32 A (ręczne, bez cewki); -K30..-K33 separacja -A3/-U2 (tryb izolacji /E1.3);\\ -K3 wentylatory --- góra arkusza /E3.8};
\end{tikzpicture}
\end{center}

\end{document}
